%! Author = sriadityavedantam
%! Date = 3/29/26

\section{Congruences}

When we talk about congruence classes mod $n$, we are essentially grouping integers based on the remainder they leave when divided by $n$.
This creates a classification system, where numbers that share the same remainder form a class.

\begin{definition}[Congruence]
    Suppose $n\nn$.
    If $a,b\nz$, we define
    \[
        a \equiv b \mod n
    \]
    as a congruence.
    We say ``$a$ is congruent to $b$ modulo $n$'' if and only if
    \[
        n \mid (b-a).
    \]
\end{definition}

\begin{lemma*}{
    $a\equiv b\mod n$ if and only if $n\mid a-b$, and this holds if and only if there exists $q\nz$ such that
    \[
        b = qn + a.
    \]
    Prove this exercise on your own.
}
\end{lemma*}

\begin{definition}[Equivalence Relation]
    Given $S$ is a set and $\sim$ is a relation on $S$, $\sim$ is an equivalence relation if for all $a,b,c\in S$,
    \begin{enumerate}
        \item $a\sim a$ (reflexive);
        \item if $a\sim b$, then $b\sim a$ (symmetric);
        \item if $a\sim b$ and $b\sim c$, then $a\sim c$ (transitive).
    \end{enumerate}
\end{definition}

We will learn and find uses for the equivalence relation, but to connect it to the topics at hand, $a\equiv b\mod n$ is an equivalence relation that envelopes congruences and what we will learn about congruence classes.
Essentially, $a\equiv b\mod n$ is the same as $a\sim b$.

\begin{lemma}{
    Congruence mod $n$ is an equivalence relation.
}
    \textbf{Case 1.}
    Let $a\nz$.
    Then $a\equiv a\mod n$, because $a-a = 0$ and $n\mid 0$.

    \textbf{Case 2.}
    Suppose $a\equiv b\mod n$.
    Then $n\mid a-b$, and due to properties of the logical divide, $n\mid b-a$.
    Thus $b\equiv a\mod n$.

    \textbf{Case 3.}
    Suppose $a,b,c\nz$, $a\equiv b\mod n$, and $b\equiv c\mod n$.
    So $n\mid b-a$ and $n\mid c-b$.
    Then $n\mid (b-a)+(c-b) = c-a$.
    Thus $a\equiv c\mod n$.
\end{lemma}

\begin{definition}[Equivalence Classes]
    Suppose $\sim$ is an equivalence relation on $S$ and $a\in S$.
    The equivalence class of $a$ is
    \[
        [a] \coloneqq \{b\in S : b\sim a\}.
    \]
\end{definition}

Let us consider an equivalence relation $\sim$ on the set of integers $\mathbb{Z}$, where $a\sim b$ if and only if $a\equiv b\mod 5$.
In this case,
\[
    [2]_5=\{\hdots,-8,-3,2,7,12,\hdots\}.
\]
This is the equivalence class of $2$, consisting of all integers that are congruent to $2$ modulo $5$.
Equivalence classes provide a systematic way of grouping elements in a set based on their relationships under an equivalence relation.

\begin{definition}[Congruence Classes]
    For a congruence mod $n$, if $a\nz$, then
    \[
        [a]\coloneqq \{b\nz : b\equiv a\mod n\}.
    \]
\end{definition}

Congruence classes provide a systematic way of grouping integers based on their remainders when divided by $n$ under a congruence relation.
They are essential in modular arithmetic, number theory, and algebraic structures, contributing to a deeper understanding of mathematical relationships and structures.Two equivalence classes are the same if they include each other.
For example, if $[a] = [b]$, then $a\in [b]$ and $b\in[a]$.
The set $S$ is the distinct union of its distinct equivalence classes.
That is, every element of $S$ is in some equivalence class.

\begin{proposition}{
    If $a,b\in S$, then either $[a] = [b]$ or $[a]\cap[b] = \varnothing$.
}
    We need to prove that if
    \[
        [a]\cap[b]\neq \varnothing,
    \]
    then
    \[
        [a]=[b].
    \]
    Let $c\in[a]\cap[b]$.
    Then $c\equiv a\mod n$ and $c\equiv b\mod n$.
    So by the next proposition,
    \[
        [c] = [a] = [b].
    \]
    Thus $[a] = [b]$.
\end{proposition}

\begin{proposition}{
    [a] = [b]\iff a\equiv b\mod n.
}
    ($\implies$)
By definition,
\[
    [a]\coloneqq\{x : x\equiv a\mod n\}.
\]
Since $a\equiv a\mod n$, we have $a\in[a]$.
If $[a]=[b]$, then $a\in[b]$.
But
\[
    [b]\coloneqq\{x\nz : x\equiv b\mod n\},
\]
so $a\equiv b\mod n$.

($\impliedby$)

\textbf{Case 1. $[a]\subseteq[b]$.}
Let $c\in[a]$.
Then $c\equiv a\mod n$.
Since $a\equiv b\mod n$, by transitivity we get $c\equiv b\mod n$.
So $c\in[b]$.
Hence $[a]\subseteq[b]$.
Similarly, we can show $[b]\subseteq[a]$.
Thus $[a] = [b]$.
\end{proposition}

This relationship provides a clear connection between the equality of equivalence classes and the congruence of integers modulo $n$.

\begin{proposition}{
    Fix $n\geq 2$.
    The distinct congruence classes modulo $n$ are
    \[
        [0],[1],\hdots,[n-1].
    \]
    In fact, if $a\nz$, then $[a]=[r]$ where $r$ is the remainder when $a$ is divided by $n$.
}
    If
    \[
        a = qn + r,\quad 0\leq r\leq n-1,
    \]
    then $a\equiv r\mod n$, so $[a]=[r]$.
    By the division algorithm, every integer $a$ has such a remainder $r$, so every class $[a]$ must be one of these classes.
    To show these classes are distinct, suppose $[i]=[j]$ where $0\leq i,j\leq n-1$.
    Then by the previous proposition,
    \[
        i\equiv j\mod n.
    \]
    So $n\mid (j-i)$.
    But since $-(n-1)\leq j-i\leq n-1$, the only multiple of $n$ in this range is $0$.
    Thus $j-i=0$, so $i=j$.
    Therefore the distinct congruence classes modulo $n$ are exactly
    \[
        [0],[1],\hdots,[n-1].
    \]
\end{proposition}